\section{Introduction}

Fluctuation relations connect nonequilibrium trajectory statistics to equilibrium free-energy differences. Jarzynski's equality and Crooks' fluctuation relation are classical results, with formulations for Markovian systems, master equations, and path ensembles \parencite{Jarzynski1997PRL,Jarzynski1997PRE,Crooks1998,Crooks1999,Crooks2000}. Continuous-time Markov chains and finite-state jump processes are likewise standard: a conservative generator determines exponential holding times, an embedded jump chain away from absorbing states, a transition semigroup $P_t=\exp(tQ)$, and explicit finite-record path densities \parencite{Norris1997,RaoTeh2013}. For a stochastic-thermodynamics reader, the payoff is that the sector construction, the $\exp(tQ)$ semigroup, the concatenated driven trajectory, and the Crooks relation are proved to refer to one and the same probability law.

What is less automatic in a proof assistant is that the objects used at different levels of the usual argument are actually the same probability law. A sector-mass calculation does not by itself identify a pushforward of a constructed path measure. A matrix exponential defined independently of the path space does not by itself describe the terminal or finite-dimensional marginals of that measure. A list of window laws does not by itself yield a measurable global real-time trajectory. Finally, an abstract detailed fluctuation identity does not by itself show that concrete forward and reverse laws generated by physical dynamics satisfy it.

This paper presents an end-to-end mechanization of that connection. The formal development begins with nonnegative jump rates on a finite state space, constructs a fixed-horizon probability measure on the disjoint union of all finite jump records, identifies the actual terminal and finite-dimensional pushforwards with the matrix-exponential semigroup, concatenates finitely many piecewise-constant windows into one real-time path, and proves physical fluctuation relations on the resulting global laws.

The main contributions are as follows.

\begin{enumerate}[leftmargin=*,label=(\arabic*)]
  \item We construct the finite-horizon jump-trajectory measure as a sum over jump-count sectors and prove its normalization, equivalently finite-horizon non-explosion.

  \item We identify the terminal law and every monotone finite-dimensional distribution of that same measure with the Markov semigroup $\exp(tQ)$.

  \item We concatenate that same measure into a stepwise driven protocol and derive the Crooks and Jarzynski relations and, for positive inverse temperature, the second-law inequality on the constructed global path space.
\end{enumerate}

The repository also contains a broader discrete-time measure-theoretic development, non-atomic Gaussian and Metropolis--Hastings examples, two-state calculations, and a refinement theorem for two discrete work conventions. Those components support and test the library, but the present paper deliberately focuses on the continuous-time chain shown in the concept figure.
\begin{figure}[htbp]
\centering
\[
\begin{aligned}
  \text{generator rates}
    &\xrightarrow{\ \text{\small free-simplex sector charts}\ }
      \text{fixed-jump-count sector laws}\\
    &\xrightarrow{\ \text{\small tail disappearance, normalization}\ }
      \Pp_{T,x}\\
    &\xrightarrow{\ \text{\small terminal and finite-time sampling}\ }
      \exp(TQ)\text{ kernels and FDDs}\\
    &\xrightarrow{\ \text{\small window kernels, measurable concatenation}\ }
      \text{forward and aligned-reverse global laws}\\
    &\xrightarrow{\ \text{\small Gibbs detailed balance}\ }
      \text{Crooks, Jarzynski, second law, work distribution}.
\end{aligned}
\]
\caption{The end-to-end theorem chain. The fixed-initial path law $\Pp_{T,x}$ is constructed from the jump-count sectors (\lean{pathLawFrom}).}
\end{figure}

\paragraph{Scope and novelty.}
The verified result concerns a finite state space, a finite total time, and finitely many stepwise changes of the energy landscape. Each dynamical window is governed by a constant, single-temperature generator that is reversible with respect to the instantaneous Gibbs distribution, and time reversal acts trivially on the states. The result does not cover smooth calendar-time dependence or a step-to-smooth refinement limit, multiple reservoirs, nonconservative steady-state forcing, odd time-reversal variables, infinite-time path laws, or diffusions. Mechanized Markov-process theory predates this work, including Isabelle/HOL developments and recent Lean infrastructure \parencite{Hoelzl2017,AFPMarkovModels,Marion2025,Ripple2026}; concurrent Lean work also studies consequences of an assumed detailed fluctuation theorem \parencite{PhysicsAIDFT}. Section~\ref{sec:related} gives the detailed positioning.
